\section{The limits of the dictionary}

\begin{statement}{G.1}{Proved}{The two fine variances have a coefficient counterpart}
The enhanced contravariant and covariant diagram theories over $S$ are
$\Fun(S^{\op},\Spaces)$ and $\Fun(S,\Spaces)$; opposition of the base
exchanges them. Their locally constant parts both identify with
$\Fun(BS,\Spaces)$. There is no inference from this last identification
to an equivalence of the two fine theories over the unchanged base, or
to Verdier duality. For an EI base their abstract stratified counterparts
are the entrance and exit variances of B.1--B.2; for a non-EI base B.7
obstructs the usual conservative stratified-space interpretation.
\end{statement}
\begin{proof}
Precomposing the first assignment with $S\mapsto S^{\op}$ gives the
second assignment literally. Inverting all arrows in either base gives
the same space, with inverse paths accounting for the variance switch.
Before that localization the two assignments remember directed arrows.
For the finite fork $S=(a\leftarrow b\to c)$, $S$ has an initial
object and no terminal object, whereas $S^{\op}$ has a terminal object
and no initial object. Both are already idempotent complete. B.6
therefore proves that their presheaf infinity-categories are not
equivalent; the same fork is checked in the program for G.2.
In a stable coefficient category, a
Verdier-duality construction would additionally require a dualizing
object and a suitable dualizable or constructible finiteness range.
Neither an operation on the indexing category nor the axioms Der~1--4
provide those data. Thus the two fine localizations discussed in VI
have a precise diagrammatic counterpart, but their interchange does not
by itself supply a six-operation duality theorem.
\end{proof}

\begin{statement}{G.2}{Proved}{Total asphericity is an exactness condition on indexing}
Let $A$ be a nonempty connected small category. The following three
conditions are equivalent: $A$ is totally $\W$-aspherical in XII
ms.~88--89; every component of the category of elements of every
covariant $H:A\to\Set$ is weakly contractible; and for every strict
$K:A\to\mathbf{sSet}$ the canonical augmentation
\[
 \hocolim_AK\longrightarrow\colim_AK
\tag{G.2}
\]
is a weak equivalence. This is a coefficient-indexing exactness property,
not a property determined by $BA$ alone and not a condition on the local
conical charts of a stratification.
\end{statement}
\begin{proof}
We verify the supplied TOTASPH B.11 reduction. The manuscript definition
says that the category of elements of $u_!1_A$ is weakly contractible for
every functor $u:A\to B$, where this Kan extension is the ordinary
presheaf Kan extension. Its value at $b$ is
$\pi_0(b\backslash A)$. The functor
\[
 A\longrightarrow B/(u_!1_A),\qquad
 a\longmapsto(u(a),[\id_{u(a)}])
\]
is ordinary final: its comma category over $(b,c)$ is exactly the
nonempty connected component $c$ of $b\backslash A$. Conversely an
ordinary final $u$ has $u_!1_A=1_B$. Totality is therefore equivalent to
every ordinary final image of $A$ being weakly contractible.

Assume totality. In particular $BA$ is contractible. Given a nonempty
connected component $E$ of $\operatorname{El}_A H$, replace $H$ by the
subfunctor consisting of that component. Adjoin a new source $t$ with
arrows $t\to a$ indexed by $H(a)$; composition is the action of $H$.
The inclusion of $A$ is ordinary final, since its only new comma
category is $E$. Its target is therefore weakly contractible. The
one-object case of C.3 identifies its classifying space with the
homotopy pushout of $BA\leftarrow BE\to *$. The reduced homology
sequence shows that $E$ is integrally acyclic.

This acyclicity property is inherited by each such $E$. Indeed for
$L:E\to\Set$ the functor
$a\mapsto\coprod_{h\in H(a)}L(a,h)$ has category of elements
isomorphic to $\operatorname{El}_E L$, by the map
$((a,h),l)\mapsto(a,(h,l))$. Fix a point of connected $BE$ and any
cyclic subgroup $C$ of its fundamental group. The coset action of
$\Pi_1(BE)$ defines a set-valued locally constant functor on $E$ whose
element category is the connected cover with fundamental group $C$.
This can be seen directly by lifting strings of arrows: a loop closes
precisely when its class lies in $C$. The inherited acyclicity makes its
first homology $C_{\mathrm{ab}}=C$ vanish. Every cyclic subgroup is
therefore trivial. Thus $BE$ is simply connected and acyclic, hence
weakly contractible by the first-nonzero-group form of Hurewicz. This
proves the second condition, including categories of elements with any
number of components. Conversely if the second condition holds, the
commas of an ordinary final $u:A\to B$ are connected categories of
elements of $a\mapsto B(b,u(a))$. They are contractible. Quillen's
Theorem A gives $BA\simeq BB$, and the constant singleton diagram makes
$BA$ contractible. This proves totality.

For a simplicial diagram $K$, its simplicial replacement has bidegrees
\[
 \coprod_{a_0\to\cdots\to a_p}K(a_0)_q.
\]
At fixed $q$ this is $N\operatorname{El}_A(K_q)$. Under the second
condition its component map to the discrete set $\colim_AK_q$ is a
weak equivalence. The bisimplicial diagonal lemma gives exactly (G.2).
Conversely apply (G.2) to a discrete set-valued diagram. Its source is
the nerve of the element category and its target its component set, so
every component must be contractible. These are the actual
augmentations, not abstract equivalences between their endpoints.

Finally let $A$ be the fork $a\leftarrow b\to c$, with its three
identities and two indicated arrows. It has initial object $b$, so
$BA$ is contractible. Set $H(b)=\{0,1\}$ and
$H(a)=H(c)=\{*\}$, with both action maps constant. Its element category
has the four objects $b_0,b_1,a,c$ and four nonidentity arrows
$b_i\to a,b_i\to c$, with no other nonidentity compositions. Its nerve
is the four-edge circle, so the second condition fails. In contrast
the point category satisfies it. The program \texttt{G\_checks.py}
checks both finite categories and their integral cycle ranks. Filtered
categories do satisfy the criterion, because finite diagrams in an
element component admit the requisite common targets and equalizations;
we make no converse assumption. In particular the exactness condition
in (G.2) is always stated when an ordinary colimit is used in place of
a homotopy colimit.
\end{proof}

\begin{statement}{G.3}{Proved}{Truncation localizers do have a sheaf meaning}
For $n\geq1$, the localizer
$W_n=\{u:\tau_{\leq n}Bu\text{ is an equivalence}\}$ is exactly the
class detected by equivalence on locally constant sheaves of
$(n-1)$-types on the CW realizations. This does not make $W_n$
equivalences stable under arbitrary ordinary homotopy pullbacks.
For every finite $n\geq0$, the map $*\to K(\mathbb Z,n+1)$ lies in
$W_n$, whereas its pullback along itself has projection
$K(\mathbb Z,n)\to *$, which is not in $W_n$.
\end{statement}
\begin{proof}
A nerve realization is locally contractible, so its locally constant
coefficients are functors on its fundamental $\infty$-groupoid. B.11
gives the stated test with $r=n-1$. For the counterexample,
$K(\mathbb Z,n+1)$ is $n$-connected, whereas the homotopy pullback of
its two basepoints is its loop space. The loop space has nonzero
$\pi_n$ if $n>0$, and is the nontrivial discrete group $\mathbb Z$ if
$n=0$. Thus the specified comparison fails at exactly the claimed
truncation level.

There is also a finite categorical $n=0$ counterexample. Let $B$ be the
one-object group $C_2$, and let $E$ be the groupoid on objects $0,1$
with one arrow $i\to j$ for every ordered pair. Define
$p:E\to B$ by sending that arrow to $j-i\pmod2$. Addition verifies
functoriality. Every base arrow has a unique lift with a fixed source
or fixed target, so $p$ is a strict bifibration. The source and target
of $p$ are connected, so $p\in W_0$. Its strict pullback along the
point of $B$ is the discrete two-object category: only arrows with
$j-i=0$ survive. Its projection to the point is not in $W_0$.
The program checks all compositions, both unique-lift properties, and
the pullback components. Thus even a strict bifibration that models
ordinary homotopy pullbacks need not preserve a larger localizer.
\end{proof}

\begin{statement}{G.4}{Known}{The intrinsic $W_n$ fibration criterion}
Let $p:T\to B$ be a functor, $n\geq0$, and define a weak
$W_n$-fibration to mean that strict pullback along $p$ preserves every
$W_n$-map of categories over $B$. This is equivalent to requiring, for
every $X\to B$, that the specified map be an equivalence of $n$-types:
\[
 \tau_{\leq n}B(T\times_BX)\longrightarrow
 \tau_{\leq n}BT
       \mathop{\times}_{\tau_{\leq n}BB}\tau_{\leq n}BX.
\tag{G.4}
\]
For $n=0$, writing $B_j$ for a connected component and
$T_j=T\times_BB_j$, this is equivalent to the two elementary conditions
that $\pi_0(T_b)\to\pi_0(T_j)$ is bijective for each $b\in B_j$, and
that each arrow $b\to b'$ has a lift inside each component of $T_j$.
\end{statement}
\begin{proof}
The properness hypothesis has a direct source. In
\citet[\S9.2.1, pp.~364--365]{E_Cisinski2006}, $W_n$ is the fundamental
localizer generated by trivializing the single category
$\Delta/\partial\Delta^{n+2}$, and it is explicitly proved proper by
Corollary~4.3.28. Corollary~9.2.15 there identifies its equivalences
with the maps inducing a bijection on $\pi_0$ and isomorphisms on
$\pi_i$ for $1\leq i\leq n$, exactly our truncation convention.
Thus its localized homotopy theory is that of $n$-types. The associated
Thomason model is right proper by Theorem~5.2.15, and its derived
pullback corresponds to the right side of (G.4), a limit of $n$-types.
Proposition~6.2.3 of \citet{E_Cisinski2006}, whose proper-localizer
hypothesis is essential, identifies weak fibrations precisely with
functors whose every strict pullback square is a derived pullback in
that theory. This proves the criterion, including the canonical
comparison map. Equivalently, factor $p$ over $B$ into a $W_n$-map
followed by a Thomason fibration. The weak-fibration comparison theorem
makes this replacement a $W_n$-map on every strict pullback; the model
pullback then computes (G.4). Conversely (G.4) sends a $W_n$-map over
$B$ to an equivalence on its right side, hence on its left side, which
is precisely the defining preservation property. These are the checks
required in the supplied POLES C.9 argument.

For clarity the $n=0$ description can be proved independently. Pull back
the component equivalence $e_b\to B_j$ to get the first condition.
Pull back the two endpoint inclusions into the arrow category
$[1]\to B$; each fibre must induce the same component bijection in
the resulting two-layer category. Components can meet its two layers
only through an arrow over the specified base arrow, giving the second
condition. Conversely fix a connected component of $X$ mapping into
$B_j$. The first condition labels every fibre component by a unique
$c\in\pi_0(T_j)$; the second condition connects equal labels across
every edge of a zigzag in $X$. Distinct labels cannot connect, since
their images in $T_j$ would then connect. This gives, naturally in $X$,
\[
 \pi_0(T\times_BX)=\pi_0(T)\times_{\pi_0(B)}\pi_0(X),
\]
including empty components. It proves the assertion at level zero
without importing an ordinary-space homotopy-pullback claim. For
$n>0$ the right side of (G.4) is a pullback \emph{of truncated spaces};
it cannot be replaced by truncation of an arbitrary ordinary pullback.
\end{proof}

\begin{statement}{G.5}{Proved}{The exact universal base-change obstruction for other localizers}
For a fundamental localizer in the relative convention of 0.1, stability
of its induced equivalences of spaces under every homotopy pullback
holds exactly for $W=\W$ and for the class of all functors. Equivalently,
every $\W$-perfect categorical functor preserves $W$ under strict
pullback only in these two cases. This is the precise obstruction; it
does not deny the coefficient meanings in B.11 or G.3--G.4.
\end{statement}
\begin{proof}
Every fundamental localizer contains $\W$, by the minimality theorem
\citep[Theorem 6.1.18]{E_Cisinski2006}. Model a space by its simplex
category, naturally and by a weak equivalence as in E.5. Model any
cospan of spaces by a strict cospan of simplicial sets, using a path
object to strictify its specified homotopy, and then by their simplex
categories. Factor one leg through the strongly proper and smooth
replacement of E.2. Such a replacement is a weak fibration for $\W$:
its strict fibres compare with commas by properness, its comma
transitions are weak equivalences by smoothness and local constancy,
and Quillen's Theorem B, after every base change, identifies its strict
pullbacks with homotopy pullbacks. This is also the application of
Cisinski's Proposition~6.2.3 in the supplied POLES E.4 proof.
Consequently categorical preservation along these perfect maps is
exactly preservation under all homotopy pullbacks of spaces.

Suppose this preservation holds and there is $f\in W\setminus\W$.
A homotopy fibre $F$ of $Bf$ is not weakly contractible. Its map to the
point lies in $W$ by pullback stability. If $F$ is empty, this already
gives $\varnothing\to *\in W$. Otherwise every choice of basepoint
$*\to F$ is in $W$, by two-out-of-three. Pulling it back along itself
gives $\Omega F\to *\in W$, and the argument can be iterated. If $F$
is disconnected, it already has two points in different components.
If it is connected, some homotopy group must be nonzero; otherwise it
would be weakly contractible. Iterating loops produces a $W$-aspherical
space $Z$ with at least two components. The two basepoint maps from
distinct components belong to $W$, but their homotopy pullback is empty.
Again $\varnothing\to *\in W$. Pulling this map back to any space $Y$
gives $\varnothing\to Y\in W$. Two-out-of-three then puts every map
$X\to Y$ in $W$. Thus $W$ is total. If no such $f$ existed, minimality
would give $W=\W$. Conversely weak homotopy equivalences, and plainly
all maps, have the stated stability.

This proof verifies POLES E.5 under the relative localizer axioms and
uses no accessibility assumption. It concerns equivalences of ordinary
spaces and the specified perfect replacement. The properness hypothesis
on the localizer in G.4 instead computes pullbacks in its localized
homotopy theory; these are different assertions, as G.3 demonstrates.
\end{proof}

\begin{statement}{G.6}{Proved}{Accessible coefficient localizations need not be geometric sheaf localizations}
Let $L:\Spaces\to\mathcal L\subset\Spaces$ be an accessible
reflective localization whose inverted maps $W$ are preserved by
product with every space. Then the constant-coefficient class of the
derivator of $\mathcal L$ is exactly $W$. Accessible fundamental
localizers admit this description through the simplex-category
correspondence. It does not follow that $L$ preserves finite limits or
that $\mathcal L$ is the sheaf $\infty$-topos of a stratified space.
\end{statement}
\begin{proof}
For $B\in\mathcal L$ and an arbitrary space $A$, the exponential
$\Map(A,B)$ is local: testing $f\in W$ against it is the same as testing
$f\times A$ against $B$. Thus mapping spaces between local objects are
local. The constant-coefficient comparison over $f:X\to Y$ is
\[
 \pi_0\Map(Y,\Map(A,B))\longrightarrow
 \pi_0\Map(X,\Map(A,B)),\qquad A,B\in\mathcal L.
\]
It is bijective whenever $f\in W$. Conversely take $A=*$, which is
local, and every local $B$. These tests say that
$\Hom_{\Ho\mathcal L}(LY,B)\to\Hom_{\Ho\mathcal L}(LX,B)$ is
bijective for every $B$. Yoneda makes $Lf$ an equivalence, so $f\in W$.
The accessible localized category is presentable and complete, giving
its derivator.

The correspondence for accessible fundamental localizers is
\citet[Theorems 4.2.15 and 1.4.3]{E_Cisinski2006}, using the test
category $\Delta$. Product stability in the categorical presentation
follows from relative Theorem A: the slices of $u\times C$ over $c$
are $u\times(C/c)$, and projection identifies them with $u$ up to
specified weak equivalences since $C/c$ has a terminal object.
This checks the product hypothesis required above and the realization
argument of GALOIS E.7. No left exactness was used. If a nontrivial
such reflector preserved finite limits, its equivalences would be
stable under all homotopy pullbacks, contrary to G.5 unless its class
were $\W$ or total. In particular finite truncation has a useful
coefficient interpretation without supplying that geometric
localization property.
\end{proof}

\begin{statement}{G.7}{Proved}{Geometric structures not supplied by the categorical data}
Neither the layered category nor the bare derivator axioms specify a
conically smooth atlas, a differentiable normal bundle, a chosen
tubular neighbourhood, a manifold dimension, or an orientation and
dualizing object. The categorical proper and smooth conditions also
do not impose a proper-support operation or a stratified submersion
structure. These additional notions therefore have no automatic
counterpart in the portions of the manuscript audited here.
\end{statement}
\begin{proof}
A one-stratum contractible manifold of any positive dimension has exit
$\infty$-groupoid equivalent to a point and the same locally constant
coefficient categories. Thus even its dimension is not determined by
that categorical datum. Smooth atlases and tubular neighbourhoods
include specified coordinate changes, derivatives and embeddings,
none of which is supplied by a functor to a poset or an indexing
derivator. An orientation or dualizing object is extra coefficient
data, as the need for a specified duality in G.1 already shows.
D.5--D.8 give completely specified finite counterexamples separating
categorical properness and smoothness from topological properness,
submersions and ambient direct image. Conically smooth stratifications
do yield useful additional structure in stratified topology, but its
presence must be stated and verified on the geometric side; it is not
reconstructed from the names \emph{propre} and \emph{lisse}.
\end{proof}
